\chapter{ОХРАНА ТРУДА И ЭКОЛОГИЧЕСКАЯ БЕЗОПАСНОСТЬ}

\section{Анализ условий труда при эксплуатации системы мониторинга}

Эксплуатация разработанного программно-аппаратного комплекса предполагает его размещение в серверных помещениях и периодическое обслуживание техническим персоналом. Основными опасными и вредными производственными факторами согласно ГОСТ 12.0.003-2015 являются:
\begin{itemize}
    \item повышенное значение напряжения в электрической цепи, замыкание которой может произойти через тело человека;
    \item повышенный уровень шума от систем охлаждения серверного оборудования;
    \item недостаточная освещенность рабочей зоны при проведении регламентных работ.
\end{itemize}

\section{Обеспечение электробезопасности}

Разработанное устройство питается от источника постоянного тока напряжением 5 В, что относится к сверхнизким напряжениям. Однако блоки питания подключаются к сети переменного тока 230 В, 50 Гц. Для обеспечения защиты персонала применяются:
\begin{enumerate}
    \item изоляция токоведущих частей (корпус устройства из ABS-пластика);
    \item применение устройств защитного отключения (УЗО) в цеппи питания серверной стойки;
    \item маркировка разъемов и кабелей в соответствии с требованиями ПУЭ.
\end{enumerate}

\section{Расчет искусственного освещения рабочего места}

Для обеспечения комфортных условий обслуживания системы произведем расчет необходимого количества светильников методом коэффициента использования светового потока. Требуемый световой поток $\Phi$ рассчитывается по формуле \eqref{eq:light}:

\begin{equation}
    \Phi = \frac{E \cdot S \cdot z \cdot K}{n \cdot \eta}
    \label{eq:light}
\end{equation}

\begin{formula_where}
    $E$ --- нормируемая освещенность (согласно СН 2.04.03-2020), лк; \\
    $S$ --- площадь помещения, $м^2$; \\
    $z$ --- коэффициент минимальной освещенности; \\
    $K$ --- коэффициент запаса; \\
    $n$ --- количество светильников; \\
    $\eta$ --- коэффициент использования светового потока.
\end{formula_where}

При норме $E = 300$ лк для помещений с ВТ, расчетное количество светильников составило 4 единицы, что обеспечивает безопасные условия для проведения монтажных работ.

\section{Экологическая безопасность и утилизация}

Проектируемое устройство не выделяет вредных веществ в процессе эксплуатации. Однако в его составе присутствуют литий-ионные аккумуляторы и печатные платы с компонентами, содержащими тяжелые металлы. 

Согласно закону РБ «Об обращении с отходами», утилизация изделия после окончания срока службы должна производиться через специализированные пункты сбора электронного лома для предотвращения попадания токсичных соединений в почву и грунтовые воды.

\newpage